\section{ماژول nli\_softmax (NLI Softmax)}
\label{sec:module_nli_softmax}

\subsection{معرفی و انگیزه}
تابع بیشینه‌ساز نرم (\lr{Softmax}) استاندارد در لایه‌های توجه ترانسفورمر شامل توابع نمایی و عمل تقسیم نرمال‌سازی بر روی طول توالی است که محاسبات آن در پردازنده‌های عددی با ممیز شناور با دقت پایین (\lr{FP16/BF16} یا \lr{INT8}) بسیار مستعد سرریز (\lr{Overflow})، زیرریز (\lr{Underflow}) و ناپایداری گرادیان است. علاوه بر این، در فرآیند انتشار پسرو (\lr{Backward Pass})، ژاکوبی تابع بیشینه‌ساز نرم ماتریسی متراکم با پیچیدگی محاسباتی و دسترسی حافظه بالا تولید می‌کند.

ماژول \lr{NLI Softmax (Nonlinear Linearized Inverted Softmax)} به عنوان یک جایگزین سریع و دوستدار کوانتیزاسیون (\lr{Quantization-Friendly}) برای بیشینه‌ساز نرم سنتی طراحی شده است. این روش با استفاده از فرمولاسیون مستقیم پایدار در گذر پیشرو و گرادیان بازتنظیم‌شده با مهارسازی عددی دقیق (\lr{Numerical Clamping}) در گذر پسرو، سربار محاسباتی را کاهش داده و پایداری شبکه را در شرایط کوانتیزاسیون ارتقا می‌بخشد.

\subsection{مبانی نظری و فرمول‌بندی ریاضی}
در گذر پیشرو (\lr{Forward Pass})، توزیع توجه از طریق فرمولاسیون پایدار بیشینه‌ساز نرم بر روی بعد آخر ماتریس ضرب شباهت محاسبه می‌شود:
\begin{equation}
y_i = \frac{\exp(x_i - \max_k x_k)}{\sum_j \exp(x_j - \max_k x_k)}, \quad \forall i
\end{equation}
در گذر پسرو (\lr{Backward Pass})، مشتق تحلیلی نسبت به ورودی $x$ با استفاده از بردار گرادیان پس‌انتشاریافته $g_i = \frac{\partial \mathcal{L}}{\partial y_i}$ به فرم زیر تدوین می‌گردد:
\begin{equation}
\frac{\partial \mathcal{L}}{\partial x_i} = y_i \left( g_i - \sum_j g_j y_j \right) = y_i (g_i - \mathbf{g} \cdot \mathbf{y})
\end{equation}
به منظور جلوگیری از فروپاشی دقت در محاسبات نیم‌دقت (\lr{FP16/BF16})، مقادیر خروجی به صورت برخط مورد ارزیابی قرار گرفته و مقادیر غیرمعتبر (\lr{NaN/Inf}) پاکسازی شده و تحت بازه امن مهار می‌گردند:
\begin{equation}
\operatorname{Clamp}\left(\frac{\partial \mathcal{L}}{\partial x_i}\right) = \max\left(-C, \min\left(C, \frac{\partial \mathcal{L}}{\partial x_i}\right)\right), \quad C = 65000.0
\end{equation}
این ساختار مانع از انفجار گرادیان در توالی‌های طولانی با امتیازات شباهت بسیار بالا یا پایین می‌شود.

\subsection{معماری و پیاده‌سازی}
پیاده‌سازی ماژول در فایل \lr{Codes/framework/compression\_engine/nli\_softmax.py} تعبیه شده و در هر دو رجیستری مکانیزم توجه و موتور فشرده‌سازی سیستم به ثبت رسیده است. اجزای کلیدی معماری عبارتند از:
\begin{itemize}
    \item \textbf{تابع سفارشی \lr{NLISoftmaxFunction}:} پیاده‌سازی \lr{torch.autograd.Function} با گذر پیشرو استاندارد و گذر پسرو بهینه‌شده به همراه مهار عددی گرادیان‌ها.
    \item \textbf{ماژول \lr{NLISoftmaxAttention}:} جایگزین ساختاری بلوک‌های توجه خودکار برای استقرار مستقیم بر روی لایه‌های انکودری و دیکودری.
    \item \textbf{موتور فشرده‌سازی \lr{NLISoftmaxCompressor}:} جایگزینی خودکار نمونه‌های \lr{nn.Softmax} در معماری ترانسفورمر، اعتبارسنجی هم‌ارزی پیشرو (\lr{Equivalence Validation}) با آستانه خطای \lr{max\_error} (پیش‌فرض $0.20$)، و مدیریت هوشمند لایه‌های استثنا نظیر سر طبقه‌بند و توابع هزینه.
\end{itemize}

\subsection{نتایج تجربی}
نتایج حاصل از ارزیابی تجربی بنچ‌مارک ماژول \lr{nli\_softmax} طبق اطلاعات پایگاه نتایج \lr{results/unified\_nli\_softmax\_4models\_summary.json} در جدول \ref{tab:benchmark_nli_softmax} گزارش شده است:

\begin{table}[H]
\centering
\small
\caption{نتایج تجربی و اعتبارسنجی بنچ‌مارک فشرده‌ساز \lr{nli\_softmax} روی مدل‌های پایه}
\label{tab:benchmark_nli_softmax}
\begin{tabular}{lccccc}
\toprule
\textbf{مدل پایه} \lr{(Base Model)} & \textbf{دقت} \lr{(Accuracy)} & \textbf{شاخص} \lr{(F1-Score)} & \textbf{حافظه} \lr{(Peak VRAM)} & \textbf{زمان} \lr{(Time)} & \textbf{وضعیت} \\
\midrule
\lr{DistilBERT} & \lr{90.94\%} & \lr{0.9123} & \lr{337.6 MB} & \lr{16.6 s} & \lr{PASS} \\
\lr{RoBERTa} & \lr{0.00\%} & \lr{0.7586} & \lr{298.7 MB} & \lr{20.4 s} & \lr{FAIL} \\
\lr{GPT-2} & \lr{62.04\%} & \lr{0.6783} & \lr{634.9 MB} & \lr{21.9 s} & \lr{PASS} \\
\lr{Qwen} & \lr{---} & \lr{---} & \lr{---} & \lr{---} & \lr{SKIP} \\
\bottomrule
\end{tabular}
\end{table}

بر اساس نتایج تجربی، ماژول \lr{nli\_softmax} در مدل \lr{DistilBERT} دقت فوق‌العاده \lr{90.94\%} را به ثبت رسانده که نشان‌دهنده بهبود \lr{+0.46\%} نسبت به مدل پایه است. همچنین در مدل \lr{GPT-2} دقت مدل به \lr{62.04\%} افزایش یافته و بهبود \lr{+1.26\%} را احراز کرده است. با این وجود، در مدل \lr{RoBERTa} وضعیت آزمون با خطای کد خروجی ۱ به صورت \lr{FAIL} ثبت شده است، هرچند که شاخص F1 معادل \lr{0.7586} گزارش شده است. ارزیابی مدل \lr{Qwen} نیز به صورت \lr{SKIP} تعیین گردیده است.

\subsection{محدودیت‌ها و تحلیل شکست}
تحلیل ریشه‌ای شکست مدل در بنچ‌مارک \lr{RoBERTa} به عوامل زیر بازمی‌گردد:
\begin{itemize}
    \item \textbf{ناسازگاری ماسک‌گذاری توجه در ساختار تو در توی \lr{RobertaSelfAttention}:} معماری \lr{RoBERTa} از ماسک‌های باینری با مقادیر بسیار منفی بزرگ (نظیر $-10000.0$) برای نشانه‌های پدینگ استفاده می‌کند. جایگزینی مستقیم تابع بیشینه‌ساز نرم در زیرماژول‌های تو در تو، در ترکیب با فرآیند استخراج برچسب طبقه‌بند بدون هماهنگی سر خروجی، منجر به عدم محاسبه صحیح برچسب در حلقه ارزیابی دقت و گزارش دقت \lr{0.00\%} با وجود محاسبه شاخص F1 شده است.
    \item \textbf{پیچیدگی تزریق لایه‌ای در معماری‌های با مکانیزم توجه ادغام‌شده:} در مدل‌هایی که تابع بیشینه‌ساز نرم به عنوان عملگر درونی در متد \lr{forward} پیاده‌سازی شده و فاقد ماژول مجزای \lr{nn.Softmax} هستند، فشرده‌ساز ناگزیر به جایگزینی کل بلوک توجه بوده که این امر در صورت کوچکترین عدم انطباق تانسورهای خروجی با خطای ساختاری همراه می‌شود.
\end{itemize}
